\begin{remark} \label{rk:auto}
Here we explain why the condition that $\phi_1, \phi_2, \phi_3$ commute in \cref{th:main-density} is not necessary if one of the $\phi_i$'s is an automorphism (see \cref{remark:density_commutative}). Without loss of generality, assume $\phi_1$ is an automorphism. Let $\psi_1 = \phi_1^{-1} \circ \phi_2$ and $\psi_2 = \textup{Id}$. Then the first two conditions of \cref{th:main-density-true} are satisfied. As for the third condition, we have $\psi_1(G) = \phi_1^{-1} \circ \phi_2(G)$ and $(\psi_1 + \psi_2)(G) = \phi_1^{-1} \circ (\phi_2 + \phi_1) (G)$. Both of these have finite index in $G$ by \cref{lem:composition}.
\end{remark}
